\section{Related Work}
\label{sec:related}

\noindent\textbf{Meshes from posed images.} MobileNeRF~\cite{mobilenerf2023}, BakedSDF~\cite{bakedsdf2023} and Binary Opacity Grids~\cite{bog2024} extract a mesh from a trained neural field and bake appearance into textures; their assets are 100--900~MB and the mesh is a by-product of a volumetric model. Meshes have also been extracted from Gaussian scenes~\cite{sugar2024,rade2024,gof2024}, at 1~GB and more, with quality below the Gaussians they came from. \MS~\cite{meshsplatting2025} and MiLo~\cite{milo2025} instead optimize the mesh itself: a triangle soup with per-vertex color and spherical-harmonic (SH) appearance, grown by splitting triangles and stitched into a connected mesh during training. These are the representations we compress; none of the above addresses the file size, and \MS's only small variant is an RGB-only PLY that drops view dependence.

\noindent\textbf{Compressing Gaussian scenes.} A large body of work reduces 3DGS~\cite{kerbl2023gaussians} from 700~MB to 5--50~MB by pruning, vector quantization and learned entropy models~\cite{compact3dgs2024,lightgaussian2024,hac2024,hacpp2025,compgs2024,rdogaussian2024,scaffold2024}. Two properties make this easier than our setting: Gaussians are independent primitives, so pruning a trained scene is cheap, and the quality loss can be fixed after training. Opaque meshes have neither property (Sec.~\ref{sec:e1}). We compare against these methods at equal file size and report the gap.

\noindent\textbf{Mesh compression.} Connectivity coders such as Edgebreaker~\cite{rossignac1999edgebreaker}, as shipped in Draco~\cite{draco}, and the MPEG SC3DMC / TFAN coder~\cite{mamou2009tfan} compress manifold meshes to about 2 bits per triangle; attributes are quantized and predicted along the traversal. We use these as baselines on our own meshes (Sec.~\ref{sec:e5}). Mesh simplification by edge collapse~\cite{garland1997qem,hoppe1996pm} is the standard way to reduce a mesh; we show it is not viable for this representation. Our growth rule follows the Lagrangian rate--distortion principle of video coding~\cite{sullivan1998rd}: a change is accepted when the distortion it removes exceeds a price times the rate it adds. The novelty is not the principle but where it is applied, inside the optimization that creates the primitives, with the rate measured by the final encoder, and a controller that turns a file size into that price.
